% !requer: abntex2-ufv.cls
% Modelo: TCC completo com a classe abntex2-ufv (variante de ../tcc-abntex2/tcc.tex)
% Compilar: latexmk -pdf tcc-ufv.tex   Testado com: abntex2-ufv 1.7, biblatex-abnt, TeX Live 2024+
% A classe (CTAN: abntex2-ufv) já define o que o modelo com a abntex2 declara à mão: os floats
% quadro, grafico e codigo com listas próprias, as listas de siglas e de abreviaturas (acro), a
% folha de aprovação e os ambientes ilustracao e tabela, que medem o conteúdo e alinham o título,
% a fonte e a nota na largura dele. Estrutura da pasta igual à do modelo com a abntex2.
%%BEGIN:ufvclasse
\documentclass{abntex2-ufv}   % 12pt, A4, frente apenas, português
%%END:ufvclasse

%%BEGIN:ufvpacotes
% 1. PACOTES: a classe carrega só o que ela precisa; estes são do trabalho
\usepackage{graphicx}
\usepackage{booktabs}
\usepackage{tabularx}
\usepackage{xcolor}
\usepackage{listings}
\usepackage{csquotes}
\usepackage[backend=biber, style=abnt, uniquename=init, giveninits,
            repeatfields]{biblatex}
\usepackage{ragged2e}
\usepackage{xurl}
\addbibresource{referencias.bib}
\appto{\bibsetup}{\RaggedRight}   % referências à esquerda (NBR 6023)
\urlstyle{same}
\lstset{basicstyle=\ttfamily\small, columns=flexible, keepspaces=true,
  numbers=left, numberstyle=\tiny, frame=single, breaklines=true,
  literate={á}{{\'a}}1 {é}{{\'e}}1 {í}{{\'i}}1 {ó}{{\'o}}1 {ú}{{\'u}}1
           {ã}{{\~a}}1 {õ}{{\~o}}1 {ç}{{\c c}}1 {ê}{{\^e}}1 {â}{{\^a}}1}
%%END:ufvpacotes

%%BEGIN:ufvsiglas
% 2. SIGLAS E ABREVIATURAS (o acro já vem com a classe)
\DeclareAcronym{abnt}{short = ABNT,
  long = Associação Brasileira de Normas Técnicas}
\DeclareAcronym{api}{short = API,
  long = Interface de Programação de Aplicações,
  long-plural-form = Interfaces de Programação de Aplicações}
\DeclareAcronym{csv}{short = CSV,
  long = \textit{Comma-Separated Values}}
\DeclareAcronym{ed}{short = ed., long = edição, tag = abreviatura}
%%END:ufvsiglas

%%BEGIN:ufvdados
% 3. DADOS DO TRABALHO (capa, folha de rosto e folha de aprovação)
\titulo{Título do trabalho}
\autor{Nome do Autor}
\local{Cidade}
\data{\the\year}
\instituicao{Nome da Instituição}
\curso{Nome do Curso}
\siglainstituicaopadrao{SIGLA}    % instituição da banca, se não informada
\orientador{Nome do Orientador}
\membrobanca{Dr. Nome do Membro da Banca}
\membrobanca[OUTRA]{Dra. Nome do Membro da Banca}
\preambulo{Trabalho apresentado à Nome da Instituição como parte das
exigências para a obtenção do título de Bacharel em Nome do Curso.}
%%END:ufvdados

\begin{document}

%%BEGIN:ufvpre
\pretextual
\imprimircapa
\imprimirfolhaderosto

\imprimirfolhadeaprovacao
% Resumo e abstract (incluídos por tcc-ufv.tex). \palavraschave e \keywords são da classe:
% separam os termos com ponto e vírgula e terminam com ponto.
\begin{resumo}
\noindent
Texto do resumo em português, em parágrafo único: objetivo, método,
resultados e conclusões do trabalho.

\palavraschave{LaTeX, trabalho acadêmico, normas ABNT}
\end{resumo}

\begin{resumo}[Abstract]
\begin{otherlanguage*}{english}
\noindent
Abstract text in English, in a single paragraph.

\keywords{LaTeX, academic writing, ABNT standards}
\end{otherlanguage*}
\end{resumo}


\listoffigures*
\cleardoublepage
\listofquadros*
\cleardoublepage
\listofcodigos*
\cleardoublepage
\listoftables*
\cleardoublepage
\listadesiglas         % só as siglas usadas no texto
\listadeabreviaturas
\tableofcontents*
\cleardoublepage
%%END:ufvpre

\textual
% Capítulo incluído por tcc-ufv.tex com \include: sem preâmbulo e sem \begin{document}.
% Na classe abntex2-ufv, \autoref de um capítulo imprime "Seção N" (o trabalho é dividido em seções).
\chapter{Introdução}\label{cap:introducao}

A elaboração de trabalhos acadêmicos segue as normas da \ac{abnt}. Este
modelo mostra as peças de um trabalho completo com a classe
\texttt{abntex2-ufv}, cada capítulo em um arquivo.

\section{Objetivos}\label{sec:objetivos}

O objetivo geral é desenvolver e avaliar um protótipo. Os objetivos
específicos são:
\begin{alineas}
  \item levantar os requisitos (\autoref{cap:metodologia});
  \item desenvolver o protótipo;
  \item avaliar o protótipo com usuários (\autoref{cap:resultados}).
\end{alineas}

\section{Organização do trabalho}

A \autoref{cap:referencial} apresenta os conceitos usados no trabalho; a
\autoref{cap:metodologia}, os métodos; a \autoref{cap:resultados}, os
resultados; e a \autoref{cap:conclusao}, as conclusões.

% Capítulo incluído por tcc-ufv.tex com \include: sem preâmbulo e sem \begin{document}.
\chapter{Referencial teórico}\label{cap:referencial}

Segundo \textcite{ex-livro}, a citação direta longa (mais de três linhas)
fica recuada:
\begin{citacao}
Texto citado literalmente, com mais de três linhas, que a classe formata
com recuo de 4\,cm, fonte menor e espaçamento simples, como orienta a NBR
10520. A referência vem ao final da citação, entre parênteses
\cite[p.~35]{ex-livro}.
\end{citacao}

Uma \ac{api} define como dois sistemas trocam dados; as \acp{api} da
\autoref{fig:arquitetura} usam arquivos \ac{csv} \cite{ex-artigo}.

% ilustracao: mede a imagem e alinha título, fonte, legenda e nota na largura dela
\begin{figure}[htbp]
\begin{ilustracao}{\caption{Arquitetura do protótipo}\label{fig:arquitetura}}%
                  {\elaboracaopropria
                   \legenda{A letra A representa o servidor.}}
  \includegraphics[width=.6\linewidth]{example-image-a}
\end{ilustracao}
\end{figure}

% Capítulo incluído por tcc-ufv.tex com \include: sem preâmbulo e sem \begin{document}.
\chapter{Metodologia}\label{cap:metodologia}

Os requisitos estão no \autoref{qua:requisitos}, e o \autoref{lst:media}
mostra a função que calcula a média das avaliações.

% quadro e codigo são floats da classe, com listas próprias
\begin{quadro}[htbp]
\begin{ilustracao}[\linewidth]{\caption{Requisitos funcionais do protótipo}%
                               \label{qua:requisitos}}%
                              {\elaboracaopropria}
  \renewcommand{\arraystretch}{1.25}
  \begin{tabularx}{\linewidth}{|c|X|c|}
    \hline
    \textbf{ID} & \textbf{Descrição} & \textbf{Prioridade} \\
    \hline
    RF01 & Cadastrar usuário com nome, \emph{e-mail} e senha. & Alta \\
    \hline
    RF02 & Exportar as avaliações em \acs{csv}. & Média \\
    \hline
  \end{tabularx}
\end{ilustracao}
\end{quadro}

\begin{codigo}[htbp]
\begin{ilustracao}[\linewidth]{\caption{Média das avaliações}\label{lst:media}}%
                              {\elaboracaopropria}
\begin{lstlisting}[language=Python]
def media(notas):
    """Devolve a média; 0 para a lista vazia."""
    return sum(notas) / len(notas) if notas else 0.0
\end{lstlisting}
\end{ilustracao}
\end{codigo}

% Capítulo incluído por tcc-ufv.tex com \include: sem preâmbulo e sem \begin{document}.
\chapter{Resultados e discussão}\label{cap:resultados}

A \autoref{tab:avaliacao} resume a avaliação, e a nota média $\bar{x}$ foi
calculada pela \autoref{eq:media}:
\begin{equation}\label{eq:media}
  \bar{x} = \frac{1}{n} \sum_{i=1}^{n} x_i
\end{equation}

% tabela: título à esquerda e conteúdo em tamanho menor (padrão IBGE), na largura da tabela
\begin{table}[htbp]
\begin{tabela}{\caption{Avaliação do protótipo por grupo de usuários}%
               \label{tab:avaliacao}}%
              {\elaboracaopropria\nota{Dados fictícios, apenas para exemplo.}}
  \begin{tabular}{lrr}
    \toprule
    \textbf{Grupo} & \textbf{Participantes} & \textbf{Nota média} \\
    \midrule
    Estudantes  & 18 & 8,4 \\
    Professores &  6 & 7,9 \\
    \bottomrule
  \end{tabular}
\end{tabela}
\end{table}

% Capítulo incluído por tcc-ufv.tex com \include: sem preâmbulo e sem \begin{document}.
\chapter{Conclusão}\label{cap:conclusao}
% Anexo: \autoref escreveria "Apêndice" (a abntex2 usa o mesmo contador);
% escreva a palavra e use \ref.

As conclusões respondem aos objetivos da \autoref{sec:objetivos}. O
questionário aplicado está no \autoref{ap:questionario}, e o regulamento
do curso (2.\,\acs{ed}), no Anexo~\ref{an:regulamento}.


\postextual
\printbibliography

\begin{apendicesenv}
\partapendices
% Apêndice (texto do próprio autor), incluído dentro de apendicesenv em tcc-ufv.tex.
\chapter{Questionário aplicado}\label{ap:questionario}

Perguntas feitas aos participantes da avaliação.

\end{apendicesenv}

\begin{anexosenv}
\partanexos
% Anexo (documento de terceiros), incluído dentro de anexosenv em tcc-ufv.tex.
\chapter{Regulamento do curso}\label{an:regulamento}

Texto do regulamento, elaborado pela coordenação do curso.

\end{anexosenv}

\end{document}
